% Figaro Paper B2: The Visibility of Coordination
% A Gramscian Reading of the Platform Economy and Its Successor
% (Political-economy half of the former Paper B; the institutional-economics
% argument is now in figaro3b1.tex.)
% Preprint, April 2026

\documentclass[11pt,a4paper]{article}

\usepackage[utf8]{inputenc}
\usepackage[T1]{fontenc}
\usepackage{amsmath,amssymb,amsthm}
\usepackage{mathtools}
\usepackage{booktabs}
\usepackage{graphicx}
\usepackage{hyperref}
\usepackage{cleveref}
\usepackage{enumitem}
\usepackage{xcolor}
\usepackage[margin=1in]{geometry}
\usepackage{natbib}
\bibliographystyle{plainnat}

\theoremstyle{remark}
\newtheorem{remark}{Remark}[section]

\newcommand{\figaro}{\textsc{Figaro}}

\title{%
  The Visibility of Coordination:\\
  A Gramscian Reading of the Platform Economy\\
  and Its Successor%
}

\author{
  Alessandro Daliana\thanks{Corresponding author. \texttt{adaliana@figaro.org}}
}

\date{April 2026}

\begin{document}

\maketitle

% ════════════════════════════════════════════════════════════════════
\begin{abstract}
\citet{gramsci1971} observed that the most durable forms of power
are those embedded so deeply in common sense that they become
invisible: the dominant arrangements rule not only through coercion
but by appearing natural, inevitable, and beyond question. This
paper applies the concept of cultural hegemony to contemporary
coordination infrastructure---the platforms, intermediaries, and
standing institutional containers through which most economic
exchange between strangers passes today.

We argue that the proposition ``two strangers need a platform to
transact safely'' has achieved hegemonic status in the precise
Gramscian sense: it is no longer argued for; it is presupposed in
the language used to describe alternatives. The same observation
applies to the propositions that production requires firms (the
Coasean assumption examined in a companion
paper~\citep{daliana2026tce}) and that cross-border trade requires
state- or alliance-controlled corridors. Each of these is a
contingent design choice that has acquired the appearance of
structural necessity.

A recently developed settlement primitive---formalized
in~\citep{daliana2026mechanism} and implemented and verified
in~\citep{daliana2026impl}---makes the coordination function
\emph{visible} by separating it from the entities that currently
perform it, and in doing so opens the organizational space to
patterns of which the firm is only one. The primitive is not a
better firm or a decentralized platform; it is an
\emph{organizational substrate} on which multiple coordination
forms can be expressed, the way electricity is a substrate on
which multiple vehicle architectures can be built. Legacy
institutions trying to describe the primitive in firm-vocabulary
are doing what internal-combustion manufacturers do when they try
to build an electric car by reusing the engine bay. This paper
is concerned not with the primitive itself but with what its
existence reveals about the incumbent arrangement: how a
contingent infrastructure passes for natural gravity, why
challenging it is qualitatively different from challenging an
explicit ideology, and what kind of communicative work is required
to make existing architecture legible as architecture.

The aim is not predictive but diagnostic. We do not argue that the
primitive will displace incumbents; we argue that recognizing such
displacement as conceivable is itself the political-economy
problem.
\end{abstract}

\medskip
\noindent\textbf{Keywords:}
cultural hegemony, Gramsci, platform economy, political economy of
infrastructure, coordination substrate, paradigm shift

% ════════════════════════════════════════════════════════════════════
\section{Introduction}
\label{sec:intro}

This paper sits between political economy and the philosophy of
infrastructure. Its starting point is a recently developed
settlement primitive that prices the cost of trust at a fixed
multiple of transaction value, with no recurring extraction and no
arbitrator. The mechanism is established in a companion
paper~\citep{daliana2026mechanism}; the implementation and formal
verification in~\citep{daliana2026impl}; the institutional-economic
implications for the boundary of the firm in a third
companion~\citep{daliana2026tce}. The present paper takes those
contributions as given and asks a different question.

If the primitive's institutional consequences are as the companion
papers suggest---if a class of standing intermediaries becomes
structurally unnecessary in a precisely defined region of
transaction space---why is the consequence so difficult to
communicate? Why does the proposal land, in venues outside its
immediate technical audience, as either obvious-but-unimportant or
implausible-by-construction? The pattern is consistent enough to
warrant explanation. The explanation we offer is Gramscian.

\paragraph{What this paper is and is not.}
This paper is an essay in the political economy of infrastructure.
It does not extend the mechanism of~\citep{daliana2026mechanism},
prove new theorems, or make empirical predictions. Its claim is
diagnostic: the difficulty of arguing for an alternative
coordination architecture is a function of the incumbent
architecture's hegemonic status, not of the alternative's intrinsic
implausibility. The argument is short because the underlying
observation is simple. We hope it is useful to those working on
infrastructure design, on platform critique, or on the
sociotechnical conditions of legitimate institutional change.

% ════════════════════════════════════════════════════════════════════
\section{Hegemony in the Gramscian Sense}
\label{sec:gramsci}

\citet{gramsci1971} distinguished two registers of class rule:
\emph{coercion} (the state's monopoly on violence) and
\emph{consent} (the cultural saturation that makes the dominant
arrangements appear inevitable). The latter, which he called
\emph{egemonia}, operates not by compelling agreement but by
constituting the field within which agreement and disagreement
become legible. A hegemonic proposition is not believed; it is
\emph{presupposed in the question}. To dispute it is not to
disagree---it is to fail to make sense.

Gramsci developed the concept in the context of class formation in
early-twentieth-century Italy, but the structure generalizes. Any
arrangement that is required to be argued for has not yet achieved
hegemony; any arrangement whose absence is unimaginable has
exceeded it. Between these poles lie the durable institutions whose
contingency is recognized in principle but functionally repressed
in practice---the institutions that one knows, on reflection, could
have been otherwise, but whose otherwise one finds oneself unable
to seriously entertain.

Gramsci's own formulation of hegemony~\citep{gramsci1971} is the
starting point; subsequent scholarship
(\citealt{anderson1976antinomies,hall1986gramsci}) has debated how
its apparatus extends beyond the class-political setting Gramsci
originally addressed. The reading we use here is closer to
\citet{hall1986gramsci}: hegemony as the saturating work that
secures consent to a particular configuration of forces by making
that configuration appear natural. The adjacent theoretical frames
we draw on are Polanyi's analysis of the disembedding of markets
from the social fabric they presuppose~\citep{polanyi1944}, the
infrastructure-studies tradition on visibility and invisibility of
coordination systems~\citep{bowker_star_1999,star1999ethnography},
and recent political-economy analyses of platform capitalism as a
particular historical form of coordination
rent~\citep{srnicek2017platform}. The present paper is a short
essay at the intersection of these traditions, not a full
engagement with any of them.

\paragraph{Three hegemonic propositions.}
We identify three propositions about contemporary coordination
infrastructure that meet the Gramscian criterion. Each is
contingent in principle; each is presupposed in practice.

\begin{enumerate}
  \item \textbf{Two strangers need a platform to transact safely.}
    The proposition is so deeply embedded that arguments
    \emph{against} platform fees are conducted in the platform's
    own register: smaller fees, fairer fees, regulated fees, public
    platforms. The proposition that no platform is required
    typically appears, in mainstream discourse, as either utopian
    or fraudulent. Yet for centuries, most economic exchange
    occurred without anything resembling the contemporary platform.
  \item \textbf{Production requires firms.}
    Examined in detail in the companion paper~\citep{daliana2026tce}:
    here we note only that the proposition is so naturalized that
    debates about ``the future of work'' are typically debates
    about the conditions of employment within firms, not debates
    about whether the firm is the right unit of production at all.
  \item \textbf{Cross-border trade requires corridors controlled by
    state or alliance actors.} Ports, railways, undersea cables,
    payment rails, settlement networks: the proposition that these
    must be controlled by some sovereign or coalition is so
    embedded that the distinction between ``geopolitical
    competition for control'' and ``geopolitical contestation of
    necessity'' is rarely made.
\end{enumerate}

These are not the only hegemonic propositions about coordination,
but they suffice for the argument. Each presents itself as
structural; each is, on inspection, contingent on a particular
historical configuration of available primitives.

\paragraph{The historical bloc.}
A Gramscian analysis is incomplete without naming the bloc that
sustains the hegemonic arrangement. The three propositions above
are sustained, jointly, by the configuration we will call the
\emph{industrial-enforcement bloc}: industrial and financial
capital that has accumulated within firm-shaped containers;
a managerial class whose authority derives from the subordination
relation that the firm institutionalizes; platform operators that
extract coordination rent at the choke points where firms,
consumers, and service providers meet; and state apparatuses whose
enforcement capacity and property-rights frameworks make the
firm-shaped container the privileged legal subject of contract,
credit, and compliance. The propositions do not float free; they
are common sense because this configuration has been the
coordination substrate for the historical period within which
modern economies have formed. A counter-hegemonic formation, in
the Gramscian sense, would require not simply an alternative
technology but the emergence of organic intellectuals whose
vocabulary makes the configuration legible as a configuration
rather than as gravity.

\paragraph{Hegemonic convergence and Schelling convergence are different.}
A distinction worth drawing: the hegemonic convergence we have
been describing is \emph{presupposed} convergence --- the
arrangement is common sense because participants do not
deliberate about it. \emph{Schelling convergence}, in the sense
of \citet{schelling1960strategy}, is \emph{recognized}
convergence --- participants knowingly pick a focal point
because they correctly expect others to pick it, and the
recognition is part of the coordination mechanism. The
contemporary platform/firm/corridor configuration operates in
the first register; the \figaro{}-aligned cohort operating
around the bonded primitive and its native token
(cf.~\citealt{daliana2026fig}) operates in the second. This
distinction sharpens the visibility argument of
\Cref{sec:coordination-design}: counter-hegemonic coordination
is necessarily a Schelling process because a counter-hegemonic
arrangement cannot be presupposed while it is still being
argued for. The work of making a coordination arrangement
legible is, in part, the work of establishing it as a credible
Schelling point around which aligned participants can converge
without central authority.

% ════════════════════════════════════════════════════════════════════
\section{The Coordination Function as Design Choice}
\label{sec:coordination-design}

The reading we propose is that the settlement primitive
of~\citep{daliana2026mechanism} makes the coordination function
visible by separating it from the entities that currently perform
it. We offer this as an interpretive claim of the present paper,
not as a theorem of the mechanism paper: Paper A establishes
equilibrium properties of a bonding game; the separation reading
is the political-economy lens through which we propose those
properties be understood. Under that lens, the platform's
coordination function (matching, enforcement, dispute resolution)
becomes a public smart contract. The firm's coordination function
(search, negotiation, enforcement) becomes three reducible costs.
The corridor's coordination function (settlement, standards,
access control) becomes a bilateral mechanism.

The crucial move is not replacement but separation. In each case,
the entity does not need to be \emph{replaced}; it needs to be
\emph{seen}. Once the coordination function is identified as a
specific bundle of services that some technology
provides---rather than as an undifferentiated property of the
entity that currently provides it---the bundle becomes
disaggregable. Disaggregation, in turn, makes substitution
conceivable.

\subsection{Figaro as Organizational Substrate}
\label{sec:substrate}

The reframing we propose is that the bonded settlement primitive is
not a better firm, a better platform, or a better corridor. It is a
\emph{substrate}: an infrastructural layer on which multiple
organizational forms can be expressed, only one of which is the
firm. Holacracy, to pick one recent example, was a network-native
coordination experiment confined to a single firm's ownership
nucleus; the primitive generalizes network-native coordination
\emph{beyond} any particular nucleus, and in doing so opens the
organizational space to patterns that the firm paradigm has no
vocabulary for --- standing treasury communities, transaction-scoped
assemblies, agent coalitions, peer networks that cross institutional
and jurisdictional lines. We take ``substrate'' here as a working
hypothesis of the present paper; the formal question of which
composed mechanisms preserve the kernel's bonding equilibrium (the
conditions under which substrate-building is defensible) is the
subject of the \emph{coordinator pattern} developed in the
implementation companion~\citep[§7]{daliana2026impl}, which we
treat as the supporting technical artifact.

The appropriate analogy is not between software tools but between
energy substrates. The transition from coal to oil and gas produced
a new industrial landscape; the internal-combustion engine was one
of its emblematic artifacts, but it was not the only one. The
current transition from fossil fuels to electric power produces a
different industrial landscape, and legacy manufacturers
repeatedly demonstrate that one cannot simply engineer the
outgoing paradigm's artifact on the new substrate and call it an
equivalent; an electric vehicle is not an internal-combustion car
with batteries attached but a distinct design whose architecture,
software, supply chain, and operational logic are native to the
new substrate.

Applied to organizational substrates: legacy institutions that try
to describe the bonded primitive in firm-vocabulary (``decentralized
firm'', ``DAO as a new kind of corporation'', ``network equivalent
of a platform'') are doing what Western automakers did for the
first two decades of the electric transition. The forms that will
dominate the new substrate are not yet named, and the vocabulary
for naming them is still being produced. This is precisely the
labor Gramsci identified: the slow, unglamorous work of producing
the conceptual vocabulary in which the new configuration becomes
evaluable on its own terms rather than as a degraded copy of the
outgoing one.

\paragraph{A note on the name.}
The protocol bears the name \figaro{} after the factotum of
Seville: the character who declares himself the city's factotum
in the famous \emph{Largo al factotum} aria of Rossini's
\emph{Il Barbiere di Siviglia} (1816, libretto by Sterbini, drawn
from Beaumarchais's \emph{Le Barbier de Séville}, 1775). Figaro is
the coordinator who arranges, brokers, and mediates between parties
of otherwise incommensurable standing without owning what is being
coordinated. He is the instrument by which commerce, marriage, and
inheritance within a noble household are actually transacted; the
household's formal principals cannot accomplish their ends without
his mediation, and the coordination he provides is not itself
legible as an institution in the classical sense. The naming is
substantive, not decorative. A bonded settlement primitive whose
construction produces a coordination function that is
enforced but not owned---that holds collateral, executes
commitments, and discharges resolution without occupying the
position of counterparty or patron---is a factotum in the literal
Renaissance sense: one that ``does everything'' while owning
nothing. The protocol is the factotum of the network on which it is
deployed; and the visibility argument we develop in this paper is,
in part, an argument about naming: an institution whose coordination
function lacks a vocabulary of its own remains structurally
invisible \emph{as} a coordination function. The protocol's name
restores one. The naming itself dates to the project's first
formulation~\citep{genovese_daliana_2022}, predating the present
mechanism by four years; the metaphor's continuity across
iterations is part of the project's intellectual record.

\begin{remark}[On What ``Visible'' Means Here]
Gramsci's analysis of hegemony emphasized the role of
\emph{intellectuals}---in his idiosyncratic sense---in producing
the conceptual vocabulary through which arrangements become
contestable. Making the coordination function visible is, in that
register, a kind of intellectual labor. It is not technical
exposition (technical readers already see it); it is the work of
producing a vocabulary in which non-technical readers can see it
without first becoming technical readers. The companion
papers~\citep{daliana2026tce, daliana2026mechanism, daliana2026impl}
do the technical work; the present paper is the prerequisite for
that work being absorbable in venues where the technical work
otherwise sounds like noise.
\end{remark}

% ════════════════════════════════════════════════════════════════════
\section{Adjacent Movements: Network States and Chambers Without Arbiters}
\label{sec:adjacent}

The Gramscian diagnostic we have developed shares family resemblance
with several contemporary movements that arrive at an overlapping
conclusion from different starting points. Three are worth naming,
both to position the present argument honestly and to identify where
the bonded primitive supplies infrastructure the movements have
otherwise struggled to specify.

\paragraph{Exit as political instrument.}
\citet{hirschman1970exit} distinguished three responses to
institutional dissatisfaction: \emph{exit} (take one's resources
elsewhere), \emph{voice} (argue within), and \emph{loyalty} (neither).
Hirschman's framework is the upstream reference for contemporary
arguments that coordination infrastructure admits of \emph{exit} as
a meaningful alternative to \emph{voice}---arguments which, absent
credible exit options, collapse back into arguments about platform
reform (voice within) or platform regulation (voice above). Our
diagnostic is, from Hirschman's perspective, a claim about the
artifactual character of the present exit constraint: the reason
exit from the platform configuration is not a serious political
option is that the coordination function has not been available
outside it. A bonded settlement primitive restores the condition
Hirschman required for exit to discipline incumbents without
recourse to voice.

\paragraph{The network-state thesis.}
\citet{srinivasan2022network} develops the Hirschmanian move for
the cloud era: communities that form online, acquire social
cohesion, build shared institutions, and eventually seek
territorial and diplomatic presence. The thesis names an
``economic engine'' as a constitutive stage between community and
recognition, but the movement has treated that stage as an
unsolved sub-problem---commerce within a network-state community in
practice continues to route through incumbent platforms and
state-clearable fiat rails, which is to say, through exactly the
hegemonic configuration the community claims to be exiting. The
bonded primitive is the missing engine. This is not a claim that
the network-state program is correct in its further aims; it is a
claim that the program requires a commerce substrate that has, at
the time of writing, no rival. A diagnostic paper is not the place
to adjudicate the normative content of adjacent movements; it is
the place to note where the present technical contribution
supplies infrastructure that adjacent movements presuppose.

\paragraph{Chambers without arbiters.}
A recurring pattern in crypto-native coordination, variously called
the decentralized autonomous chamber of commerce, the coordinator
DAO, or under earlier branding the decentralized autonomous
company~\citep{larimer2014dac}, is the attempt to federate merchant
or participant coordination without incorporating a governance body
empowered to adjudicate disputes. The attempts recurrently drift
into governance because participants have no other tool for the
coordination function: once dispute resolution is on the table, a
voting apparatus is the default answer. The bonded primitive
provides the tool. Asymmetric bonding handles the enforcement
function that the chamber would otherwise have to adjudicate;
permissionless operator registration handles the membership
function that the chamber would otherwise have to gatekeep;
transaction-scoped assemblies handle the coordination function
that the chamber would otherwise have to govern. The
chamber-without-arbiters ambition is not obviated by the primitive;
it is made architecturally achievable, which it was not before.

\paragraph{Climate and cross-sovereignty coordination.}
Climate change is a coordination problem that has resisted
solution in part because its scale exceeds any single sovereignty.
The voluntary carbon market, established to coordinate
emissions-reduction credit across jurisdictional boundaries, has
been wracked by integrity
crises~\citep{west2020overcrediting}: many credits bear little
relationship to the emissions reductions they claim. The
structural failure is that buyers cannot verify credits directly
and have no choice but to trust a registry. The bonded primitive
offers a different arrangement. A buyer bonded with a
carbon-removal operator can require schema-validated GHG
attestations~\citep[§4]{daliana2026impl}, verify the attestations
against the primitive's schema registry, and hold the operator's
bond until resolution. This is not a claim that \figaro{} solves
climate change; it is a claim that the infrastructural failure
mode of the voluntary carbon market --- trusted registries
producing credits without verifiable provenance --- is of the same
structural type as the coordination-function failure mode of the
platform economy this paper has been analyzing, and that the
primitive addresses both in the same way.

\paragraph{Position of the present paper.}
The Gramscian framing is neither a substitute for nor a rival to
these movements. It supplies the diagnostic register within which
the coordination-infrastructure question can be posed at all; the
movements supply concrete political and organizational programs
within which the bonded primitive can be deployed. A reader
arriving from the network-state literature should read this paper
as naming the hegemonic configuration their program is trying to
exit; a reader arriving from the coordinator-DAO literature should
read it as naming the reason their coordinator-without-governance
attempts have been structurally difficult.

% ════════════════════════════════════════════════════════════════════
\section{Why This Argument Is Hard to Make}
\label{sec:hardness}

The practical difficulty of communicating the implications of
infrastructure shifts of this kind is not, we argue, a
communicative deficit on the part of the proponent. It is a
structural feature of arguing against an arrangement whose
necessity has achieved hegemonic status.

\paragraph{The asymmetry of argument.}
Proposing to remove a structure that most people do not know is
there is qualitatively different from proposing a better version
of the same structure. The latter operates within the existing
register of debate (more efficient platforms, fairer firms, more
inclusive corridors); the former requires first establishing the
existence of the register itself. The proponent is not asking the
listener to accept a counter-argument; they are asking the
listener to accept that there is something to argue about. This is
much harder.

\paragraph{The evaluator's predicament.}
A skeptical reader evaluating an architectural proposal that
implies the dissolution of a major incumbent must do two
unfamiliar things at once: (a) accept that the incumbent's
necessity is contingent, and (b) evaluate a specific alternative.
Either move is hard alone; together they are uncomfortable enough
that many readers, in good faith, decline both and conclude that
the proposal is implausible. The conclusion is rational under the
listener's actual evaluative budget; it is not evidence about the
proposal.

\paragraph{The proponent's inversion.}
A proponent of an architectural alternative typically frames the
work in terms of the incumbent: ``A better Uber'', ``Uber without
the middleman'', ``a decentralized firm''. This framing has two
costs. First, it concedes the centrality of the incumbent: the
alternative is defined relative to it. Second, it re-encodes the
incumbent's hegemonic vocabulary into the alternative, importing
the very assumptions the alternative is supposed to dissolve. We
argue that this is structurally a mistake. The right framing is
disjunctive: the alternative does what the incumbent does
\emph{plus} what the incumbent prevents, where the latter is the
substantive contribution. ``A better X'' frames are
near-guaranteed to be misread; the work of finding non-incumbent
vocabulary is the political-economy work the technology cannot do
for itself.

\paragraph{The role of repeated exposure.}
Hegemonic propositions are dissolved not by definitive argument
but by sustained exposure to alternative arrangements that
function. The role of working examples (process trees that resolve
without a platform, treasuries that coordinate without a firm,
corridors that settle without a sovereign) is not to prove that
the alternative is possible; it is to gradually shift the
horizon of the imaginable. This is the slow, not the dramatic,
register of institutional change.

% ════════════════════════════════════════════════════════════════════
\section{Conclusion}
\label{sec:conclusion}

This paper has argued that the platform, the firm, and the
corridor share a structural property: each is a contingent
coordination architecture that has achieved hegemonic status in
the Gramscian sense. The recently developed settlement primitive
of~\citep{daliana2026mechanism} makes the coordination function
performed by each entity visible as a separable design artifact.
This visibility is the prerequisite for any subsequent
substitution argument; it does not, by itself, settle whether
substitution will occur, at what rate, or in which regions of the
transaction space.

The political-economy stakes are real but should not be
overstated. The argument is diagnostic rather than predictive: a
proposition that is presupposed cannot be evaluated; making the
proposition explicit is the work that political economy must do
before institutional analysis can proceed. A more cautious version
of the same point: in a hegemonic arrangement, the substantive
work of an architectural proposal is largely the work of producing
the vocabulary in which the proposal becomes evaluable. The
technology rarely does this work for itself.

We close with two observations of method. First, the kernel is
ideologically agnostic; the graph is the politics. The
irreducible settlement
primitive~\citep[§1.4]{daliana2026mechanism}---the
\texttt{commit}-and-\texttt{resolveProcess} operations, the bond
structure, and the equilibrium they induce---contains no
commitments about currency, jurisdiction, identity, arbitration,
role structure, price-discovery, or contribution metric. All of
these are expressed in the graph that parties compose on top of
it: a Dutch-auction-and-private-operator graph expresses a
market-liberal pattern; a cooperative-bonding-and-contribution-weighted-resolution
graph expresses a cooperative one; an interest-free risk-sharing
graph expresses an Islamic-finance one. The kernel enables each
and prefers none. This is not rhetorical evasion but the
architectural consequence of placing the enforcement function at
TCP/IP altitude: a layer below which nothing remains to factor
out, and a layer above which normative choices become legible as
choices rather than as properties of the infrastructure. A
defender of platforms can use the visibility move of this paper
to clarify what is actually being defended; a critic can use it
to clarify what is actually being criticized; neither has to
contest the primitive itself. Second, the hegemony framework
applies recursively. If the bonded settlement
primitive itself becomes hegemonic in turn---if its necessity
becomes presupposed in the way that platform necessity is
presupposed today---the same diagnostic work will be required to
make \emph{that} contingency visible. There is no terminal
infrastructure; there is only the next infrastructure whose
contingency we have not yet learned to see.

\section*{Acknowledgments}
I thank the readers who pressed me to separate this argument from
the institutional-economics argument with which it had been
entangled in earlier drafts. The two address different audiences
and demand different evidentiary standards; mixing them did
neither audience justice. The Gramscian framing also owes a debt
to several conversations with collaborators inside and adjacent to
the Coalition of Automated Legal Applications (COALA), whose work
on the political economy of automated legal infrastructure has
shaped how I think about the question of visibility.


% ════════════════════════════════════════════════════════════════════
% References
% ════════════════════════════════════════════════════════════════════

\begin{thebibliography}{99}

\bibitem[Daliana(2026a)]{daliana2026mechanism}
Daliana, A.
\newblock Asymmetric Bonding: A Self-Enforcing Settlement Primitive
  for $N$-Party Coordination.
\newblock Companion mechanism-design paper, 2026.

\bibitem[Daliana(2026b)]{daliana2026impl}
Daliana, A.
\newblock The \textsc{Figaro} Kernel: A Verified Solidity
  Implementation of Asymmetric Bonding.
\newblock Companion implementation paper, 2026.

\bibitem[Daliana(2026c)]{daliana2026tce}
Daliana, A.
\newblock From Firms to Transaction-Scoped Institutions: A Coasean
  Re-Examination Under Costless Bilateral Enforcement.
\newblock Companion institutional-economics paper, 2026.

\bibitem[Daliana(2026d)]{daliana2026fig}
Daliana, A.
\newblock \textsc{Fig}: A Schelling-Point Token for the
  \textsc{Figaro} Coordination Ecosystem.
\newblock Companion cryptoeconomics paper, 2026.

\bibitem[Anderson(1976)]{anderson1976antinomies}
Anderson, P.
\newblock The Antinomies of Antonio Gramsci.
\newblock \emph{New Left Review}, I/100:5--78, 1976.

\bibitem[Bowker and Star(1999)]{bowker_star_1999}
Bowker, G.~C. and Star, S.~L.
\newblock \emph{Sorting Things Out: Classification and Its
  Consequences}.
\newblock MIT Press, Cambridge, MA, 1999.

\bibitem[Genovese and Daliana(2022)]{genovese_daliana_2022}
Genovese, F.~R. and Daliana, A.
\newblock Figaro: A Proof-of-Delivery Blockchain Solution to the
  Last Mile Problem.
\newblock Unpublished working paper, March 2022.
\newblock Archived in the project repository at
  \texttt{docs/archive/origins/}.

\bibitem[Gramsci(1971)]{gramsci1971}
Gramsci, A.
\newblock \emph{Selections from the Prison Notebooks}.
\newblock International Publishers, New York, 1971.
\newblock Edited and translated by Q.~Hoare and G.~Nowell-Smith.

\bibitem[Hall(1986)]{hall1986gramsci}
Hall, S.
\newblock Gramsci's Relevance for the Study of Race and Ethnicity.
\newblock \emph{Journal of Communication Inquiry}, 10(2):5--27, 1986.

\bibitem[Hirschman(1970)]{hirschman1970exit}
Hirschman, A.~O.
\newblock \emph{Exit, Voice, and Loyalty: Responses to Decline in
  Firms, Organizations, and States}.
\newblock Harvard University Press, Cambridge, MA, 1970.

\bibitem[Larimer(2014)]{larimer2014dac}
Larimer, D.
\newblock Overpaying for Security: The Hidden Costs of Bitcoin.
\newblock \emph{Let's Talk Bitcoin}, September 2014.
\newblock Early articulation of the decentralized autonomous
  company (DAC) framing.

\bibitem[Srinivasan(2022)]{srinivasan2022network}
Srinivasan, B.
\newblock \emph{The Network State: How to Start a New Country}.
\newblock Self-published, 2022.

\bibitem[Polanyi(1944)]{polanyi1944}
Polanyi, K.
\newblock \emph{The Great Transformation: The Political and
  Economic Origins of Our Time}.
\newblock Farrar \& Rinehart, New York, 1944.

\bibitem[Schelling(1960)]{schelling1960strategy}
Schelling, T.~C.
\newblock \emph{The Strategy of Conflict}.
\newblock Harvard University Press, Cambridge, MA, 1960.

\bibitem[Srnicek(2017)]{srnicek2017platform}
Srnicek, N.
\newblock \emph{Platform Capitalism}.
\newblock Polity, Cambridge, 2017.

\bibitem[Star(1999)]{star1999ethnography}
Star, S.~L.
\newblock The Ethnography of Infrastructure.
\newblock \emph{American Behavioral Scientist}, 43(3):377--391, 1999.

\bibitem[West et al.(2020)]{west2020overcrediting}
West, T.~A.~P., B\"orner, J., Sills, E.~O., and Kontoleon, A.
\newblock Overstated Carbon Emission Reductions from Voluntary
  REDD+ Projects in the Brazilian Amazon.
\newblock \emph{Proceedings of the National Academy of Sciences},
  117(39):24188--24194, 2020.

\end{thebibliography}

\end{document}
